The conclusions of this replication split in two. Firstly, the mechanism transfers, since
the AGOP recovers linear concept directions even at the scale of 8.73M parameters,
reaching a validation $|\rho|$ as high as 0.978, and those directions are causal when
written back into the residual stream. Secondly, the capability that the steering is
supposed to bring does not appear, given that at this scale the intervention moves the
vocabulary the model reaches for long before it produces a coherent sentence carrying the
concept, so that we obtain a mechanism without the behaviour it was meant to support.

The comparative claim, in turn, never received a fair test. RFM does beat PCA cleanly, at
a validation $|\rho|$ of 0.918 against 0.154 on \verb|valence| and a monitoring AUROC of
1.000 against an overall 0.879, and yet RFM, difference-in-means and logistic regression
sit together at the AUROC ceiling, with logistic regression edging RFM out on validation
$|\rho|$ on all three concepts. Note that this says the three methods are indistinguishable
\emph{by these metrics on these concepts}, and it does not say that the AGOP recovers the
same direction as a subtraction of two class means, which is a separate claim, testable
through the per-block alignment
$\cos\left(v_{\mathrm{RFM}}, v_{\mathrm{DiffMeans}}\right)$, and one we did not compute.

One methodological consequence of this work is portable to any scale, and it is not a
result about GuppyLM. A steering coefficient must be reported relative to the activation
scale, in the form
$\varepsilon_\ell = r \cdot \operatorname*{median}_{i} \left\| a_\ell^{(i)} \right\|_2$,
since the per-block norms measured here run from 11.9 to 30.9 while the absolute values of
Beaglehole et al. are calibrated to Llama's activation scale, so that it is the relative
coefficient $r$, and not $\varepsilon$, that ports across models. Note that a replication
carried out at any other scale inherits this problem before it inherits any other.

To sustain the further substantive claims that Beaglehole et al. (\cite{beaglehole2026toward})
make, namely the scaling of steerability with model size, the advantage of RFM on hard
concepts, and the real monitoring benchmarks such as FAVABENCH, HaluEval and ToxicChat, we
require infrastructure that a single GuppyLM checkpoint cannot provide. We shall therefore
take a Pythia model ladder as a follow-on plan, analysing checkpoints across several model
sizes so as to read steerability directly against scale rather than inferring it. The main
question there is whether the ladder begins to produce fluent text carrying the concept
rather than vocabulary drift, and at which size it begins to do so.

In conclusion, the machinery of Beaglehole et al. appears not to require a large model in
order to function, but only in order to be worth using, and it is the distance between
those two statements that a single tiny checkpoint is able to measure. Here we have
discussed one model, one family of in-domain concepts and two imperfect judges, and much
more is yet to be discussed, particularly the behaviour of the same pipeline on concepts
that a model of this size does not already state lexically in its own output.
